\section{Methodology}
\label{sec:methodology}

\subsection{Baseline Framework}

We build on the LTR scheduling framework of~\cite{fu2024efficient}, which
trains an OPT-125M predictor with a ListMLE ranking loss to score prompts by
expected relative output length and uses that score to reorder the admission
queue before requests reach vLLM's generation engine. This ranking step
operates entirely outside vLLM's own continuous-batching scheduler, which
handles the resulting iteration-level batching and memory management
unmodified, as shown in Figure~\ref{fig:architecture_pipeline}.

\subsection{Proposed Extension: Composite Scoring}

Table~\ref{tab:scheduling_policies} summarizes the four scheduling policies
compared in this work. For the proposed extension, each request's admission
priority is
\begin{equation}
    \mathrm{score}(x) = \alpha \, s_{\mathrm{LTR}}(x) \;-\; \beta \, \frac{\mathrm{SCALE}}{\mathrm{len}(x)}
    \label{eq:composite-score}
\end{equation}
where $s_{\mathrm{LTR}}(x)$ is the predictor's learned score for prompt $x$,
$\mathrm{len}(x)$ is its tokenized length, $\mathrm{SCALE}$ brings the length
term onto a comparable numeric range to $s_{\mathrm{LTR}}(x)$, and
$\alpha,\beta$ are fixed weights with $\alpha+\beta=1$. The length term is
subtracted so that longer prompts, absent strong evidence to the contrary
from $s_{\mathrm{LTR}}$, receive a comparatively lower score, reflecting the
assumption that longer prompts more often elicit longer completions.
Requests are admitted in decreasing order of $\mathrm{score}(x)$; since
$\mathrm{len}(x)$ comes directly from tokenization, this signal adds no
inference cost beyond the forward pass already needed for $s_{\mathrm{LTR}}(x)$.

The predictor used for $s_{\mathrm{LTR}}(x)$ in the extension is retrained
with a pairwise margin ranking loss rather than ListMLE, so its score
distribution suits combination with the length term; both share the same
OPT-125M backbone and training set. Table~\ref{tab:hyperparameters} lists the
composite-scoring weights and training hyperparameters used throughout
Section~\ref{sec:evaluation}.

\begin{table}[t]
    \centering
    \caption{Scheduling policies compared in this work.}
    \label{tab:scheduling_policies}
    \begin{tabular}{@{}p{1.5cm}p{5.4cm}@{}}
        \toprule
        Policy & Admission ordering rule \\
        \midrule
        FCFS & Requests admitted in arrival order; no reordering. \\
        Classification & Requests bucketed into coarse predicted output-length classes; shorter buckets prioritized. \\
        LTR & Requests ordered by the OPT-125M predictor's learned ranking score, $s_{\mathrm{LTR}}(x)$, alone. \\
        LTR + Prompt Length (ours) & Requests ordered by the composite score in Eq.~\eqref{eq:composite-score}, combining $s_{\mathrm{LTR}}(x)$ with tokenized prompt length. \\
        \bottomrule
    \end{tabular}
\end{table}

\begin{table}[t]
    \centering
    \caption{Composite scoring and predictor training hyperparameters.}
    \label{tab:hyperparameters}
    \begin{tabular}{@{}ll@{}}
        \toprule
        Parameter & Value \\
        \midrule
        $\alpha$ (LTR score weight) & 0.7 \\
        $\beta$ (prompt-length weight) & 0.3 \\
        SCALE (length-term scale factor) & 100 \\
        Predictor backbone & OPT-125M \\
        Original predictor loss & ListMLE \\
        Extension predictor loss & Pairwise margin \\
        Margin ($m$) & 1.0 \\
        Pair filter threshold ($\delta$) & 0.20 \\
        Training epochs & 10 \\
        \bottomrule
    \end{tabular}
\end{table}


\subsection{Implementation Correction}
\label{subsec:sign-correction}

An early implementation of Eq.~\eqref{eq:composite-score} added the length
term rather than subtracting it, inverting its intended effect. This was
corrected prior to the experiments in Section~\ref{sec:evaluation}; we note
it here for transparency relative to any earlier, uncorrected version of this
codebase.
